\documentclass[11pt,letterpaper]{article}
\usepackage{cornell-notes}
\usepackage{needspace}

\title{Chapter 9: UNIX I: Privileges and Files}
\author{}
\date{}

\setDocTitle{Chapter 9: UNIX I: Privileges and Files}
\setDocSubtitle{Application Security Review}
\setCornellCollection{Software Security Assessment}
\setCornellUnitType{Chapter}
\setCornellUnitNumber{9}
\setCornellUnitTitle{UNIX I: Privileges and Files}
\setCornellCompanionLabel{Study and audit reference}
\setCornellSourceRange{pp. 459--557}
\begin{document}
\maketitle

\begin{center}
\small\color{Meta}\textbf{Coverage range:} pp. 459--557 \quad\textbullet\quad \textbf{Format:} Cornell cue-and-notes
\end{center}
\vspace{0.25em}
\begin{CornellOverviewBox}{Review Orientation}
\textbf{Central problem.} Audit UNIX identity, privilege transitions, path resolution, file permissions, links, races, temporary files, and stream I/O as one coherent authority model.
\textbf{Why it matters.} Security reviewers use this material to connect attacker-controlled conditions to violated invariants, consequential operations, and defensible remediation.
\textbf{Prerequisites.} UNIX users and groups, processes, filesystem namespaces, permissions, links, descriptors, and C file APIs.
\textbf{Assessment relationship.} It applies trust and lifetime analysis to operating-system resources that frequently mediate privileged application behavior.
\end{CornellOverviewBox}
\section{Learning Objectives}
\begin{itemize}
\item Explain the security significance of users, groups, and processes.
\item Identify attacker influence and failed assumptions associated with privileged programs.
\item Trace security-relevant data or state through permanent privilege drop.
\item Evaluate whether controls addressing temporary privilege drop preserve the intended invariant.
\item Audit implementation and error paths related to file and directory permissions.
\item Distinguish safe and unsafe uses of file creation.
\item Diagnose a failure involving directory safety and dangerous places from concrete evidence.
\item Verify remediation and test coverage for filenames and paths.
\end{itemize}
\section{Cornell Cue-and-Notes}
\Needspace{8\baselineskip}
\subsection{UNIX Identity and Privilege}
\begin{longtable}{>{\raggedright\arraybackslash\columncolor{CornellCueBackground}}p{0.285\textwidth}|>{\raggedright\arraybackslash}p{0.625\textwidth}}
\rowcolor{Navy}\color{white}\textbf{Cue / Recall Prompt} & \color{white}\textbf{Technical Notes and Audit Reasoning} \\
\endfirsthead
\rowcolor{Navy}\color{white}\textbf{Cue / Recall Prompt} & \color{white}\textbf{Technical Notes and Audit Reasoning} \\
\endhead
\term{Users, groups, and processes}\\[-0.1em]{\small Which identity does the kernel use for an access decision?} & Processes carry multiple user and group identities, and files carry owners and mode bits. The relevant identity can differ across system calls and privilege transitions.\par\smallskip\textbf{Audit move:} Record real, effective, saved, and supplementary identities at each sensitive operation. \\\hline
\term{Privileged programs}\\[-0.1em]{\small Why are set-user-ID and similar programs high risk?} & They process less-trusted input while holding another principal's authority. Environment, files, descriptors, signals, and invoked programs become part of the attack surface.\par\smallskip\textbf{Audit move:} Minimize privileged code and inventory every input accepted before privilege is reduced. \\\hline
\term{Permanent privilege drop}\\[-0.1em]{\small What proves that privilege cannot be regained?} & A permanent drop must update all relevant identities and groups, check errors, and verify the resulting state. Partial changes may leave a saved identity or capability that restores authority.\par\smallskip\textbf{Audit move:} Check every identity-changing return value and test attempted regain. \\\hline
\term{Temporary privilege drop}\\[-0.1em]{\small What additional state must be controlled?} & Temporarily changing effective identity can support a narrow task, but restoration, signals, threads, descriptors, and concurrent work complicate the model. Untrusted work must not run while elevated.\par\smallskip\textbf{Audit move:} Bracket the minimal operation and verify state on normal and exceptional exits. \\\hline
\end{longtable}
\begin{CornellSummaryBox}{Section Summary}
Real, effective, and saved identities determine what a process can access and whether dropped authority can return.
\end{CornellSummaryBox}
\Needspace{8\baselineskip}
\subsection{File Security and Paths}
\begin{longtable}{>{\raggedright\arraybackslash\columncolor{CornellCueBackground}}p{0.285\textwidth}|>{\raggedright\arraybackslash}p{0.625\textwidth}}
\rowcolor{Navy}\color{white}\textbf{Cue / Recall Prompt} & \color{white}\textbf{Technical Notes and Audit Reasoning} \\
\endfirsthead
\rowcolor{Navy}\color{white}\textbf{Cue / Recall Prompt} & \color{white}\textbf{Technical Notes and Audit Reasoning} \\
\endhead
\term{File and directory permissions}\\[-0.1em]{\small What authority follows from directory permissions?} & Directory search, create, remove, and rename rights differ from permissions on the target file. A protected file can still be replaced or redirected when an attacker controls a containing directory.\par\smallskip\textbf{Audit move:} Audit permissions and ownership on the full path, not only the final object. \\\hline
\term{File creation}\\[-0.1em]{\small How should a new file acquire secure permissions?} & Creation mode, process mask, inherited context, flags, and preexisting names determine the resulting object. Opening and then fixing permissions leaves a window and may operate on an attacker-chosen file.\par\smallskip\textbf{Audit move:} Use atomic creation with restrictive initial permissions and verify the opened object. \\\hline
\term{Directory safety and dangerous places}\\[-0.1em]{\small Which directories are unsuitable for privileged path operations?} & World-writable or attacker-controlled directories permit name replacement, link insertion, and cleanup races. Relative paths also depend on mutable working-directory state.\par\smallskip\textbf{Audit move:} Avoid privileged resolution through untrusted directories and use trusted directory handles where possible. \\\hline
\term{Filenames and paths}\\[-0.1em]{\small Why is textual path validation insufficient?} & Repeated separators, parent components, alternate roots, mount changes, and links can make textual form diverge from the resolved object. The namespace can change between operations.\par\smallskip\textbf{Audit move:} Bind policy to a resolved object or trusted directory-relative operation rather than a string alone. \\\hline
\end{longtable}
\begin{CornellSummaryBox}{Section Summary}
Path-based access is a name-resolution process influenced by every directory component, link, mount, and caller-controlled string.
\end{CornellSummaryBox}
\Needspace{8\baselineskip}
\subsection{File Internals and Links}
\begin{longtable}{>{\raggedright\arraybackslash\columncolor{CornellCueBackground}}p{0.285\textwidth}|>{\raggedright\arraybackslash}p{0.625\textwidth}}
\rowcolor{Navy}\color{white}\textbf{Cue / Recall Prompt} & \color{white}\textbf{Technical Notes and Audit Reasoning} \\
\endfirsthead
\rowcolor{Navy}\color{white}\textbf{Cue / Recall Prompt} & \color{white}\textbf{Technical Notes and Audit Reasoning} \\
\endhead
\term{Descriptors, inodes, and directories}\\[-0.1em]{\small What does an open descriptor guarantee?} & A descriptor refers to an opened object independently of later renaming, while an inode identifies filesystem metadata within its context. Descriptor reuse and inherited descriptors can still cause confusion.\par\smallskip\textbf{Audit move:} Verify type, owner, permissions, and flags on the opened descriptor before use. \\\hline
\term{Symbolic links}\\[-0.1em]{\small How do symbolic links redirect privileged access?} & A symbolic link stores another path and can redirect lookup across directories or filesystems. Attackers exploit applications that check one pathname but later follow a substituted link.\par\smallskip\textbf{Audit move:} Use no-follow and directory-relative primitives where supported, then verify the opened object. \\\hline
\term{Hard links}\\[-0.1em]{\small Why can a second name affect policy?} & Hard links give one underlying object multiple names. Name-based ownership or location assumptions may fail when an attacker can create a link to a sensitive object.\par\smallskip\textbf{Audit move:} Review link-creation permissions and avoid treating pathname uniqueness as object identity. \\\hline
\end{longtable}
\begin{CornellSummaryBox}{Section Summary}
Descriptors and object identifiers are more stable than names, but reuse and link semantics still require careful verification.
\end{CornellSummaryBox}
\Needspace{8\baselineskip}
\subsection{Race Conditions and Temporary Files}
\begin{longtable}{>{\raggedright\arraybackslash\columncolor{CornellCueBackground}}p{0.285\textwidth}|>{\raggedright\arraybackslash}p{0.625\textwidth}}
\rowcolor{Navy}\color{white}\textbf{Cue / Recall Prompt} & \color{white}\textbf{Technical Notes and Audit Reasoning} \\
\endfirsthead
\rowcolor{Navy}\color{white}\textbf{Cue / Recall Prompt} & \color{white}\textbf{Technical Notes and Audit Reasoning} \\
\endhead
\term{TOCTOU}\\[-0.1em]{\small What changes between check and use?} & A pathname check followed by a separate open, create, chmod, or delete creates a race if another principal can replace a component or target. Faster checks do not make the sequence atomic.\par\smallskip\textbf{Audit move:} Replace check-then-use sequences with one atomic operation or descriptor-based validation. \\\hline
\term{Permission and ownership races}\\[-0.1em]{\small Can access rights change mid-operation?} & Multiple calls that establish ownership, mode, and content may expose intermediate states. Signals, threads, and cooperating processes can enlarge the race window.\par\smallskip\textbf{Audit move:} Create securely from the start and verify postconditions on the same descriptor. \\\hline
\term{Unique temporary files}\\[-0.1em]{\small What properties must a temporary file have?} & The name must be unpredictable or atomically exclusive, the directory trustworthy, permissions restrictive, and cleanup bound to the correct object. Predictable-name generation followed by open is unsafe.\par\smallskip\textbf{Audit move:} Use a secure creation primitive and retain the returned descriptor. \\\hline
\term{Reuse and cleaners}\\[-0.1em]{\small How can cleanup create a vulnerability?} & Reusing names or relying on external cleaners can remove, replace, or expose a file while the application still expects it. Long-lived temporary objects invite namespace races.\par\smallskip\textbf{Audit move:} Unlink early when appropriate and coordinate lifetime without trusting pathname persistence. \\\hline
\end{longtable}
\begin{CornellSummaryBox}{Section Summary}
Separate checks and uses are unsafe when an attacker can mutate the namespace or permissions between them.
\end{CornellSummaryBox}
\Needspace{8\baselineskip}
\subsection{Stdio Interface}
\begin{longtable}{>{\raggedright\arraybackslash\columncolor{CornellCueBackground}}p{0.285\textwidth}|>{\raggedright\arraybackslash}p{0.625\textwidth}}
\rowcolor{Navy}\color{white}\textbf{Cue / Recall Prompt} & \color{white}\textbf{Technical Notes and Audit Reasoning} \\
\endfirsthead
\rowcolor{Navy}\color{white}\textbf{Cue / Recall Prompt} & \color{white}\textbf{Technical Notes and Audit Reasoning} \\
\endhead
\term{Opening streams}\\[-0.1em]{\small Which mode and pathname assumptions affect stream safety?} & Stream modes can create, truncate, append, or reuse files, and path lookup retains the same race concerns as lower-level open operations. Conversion from descriptors must preserve flags and ownership.\par\smallskip\textbf{Audit move:} Review mode strings, creation semantics, and the security of the resolved object. \\\hline
\term{Reading and writing}\\[-0.1em]{\small Can partial or buffered I/O change program state?} & Reads and writes may process fewer elements than requested, encounter end-of-file, or defer errors until a flush. Text transformations and formatted conversion add parsing risks.\par\smallskip\textbf{Audit move:} Check counts, error indicators, buffer state, and the distinction between EOF and failure. \\\hline
\term{Closing}\\[-0.1em]{\small Why does close handling matter?} & Buffered data may fail during flush or close, so earlier successful writes do not prove durable output. Cleanup errors can also interact with retries and descriptor reuse.\par\smallskip\textbf{Audit move:} Check finalization errors and avoid retry logic that can act on a reused descriptor. \\\hline
\end{longtable}
\begin{CornellSummaryBox}{Section Summary}
Buffered streams add state and error semantics above descriptors that must be handled explicitly.
\end{CornellSummaryBox}
\begin{CornellLegacyBox}{Historical Context}
Available descriptor-relative and no-follow interfaces differ across UNIX variants and versions. Assess the exact target platform and prefer its current race-resistant primitives.
\end{CornellLegacyBox}
\section{Security-Review Checklist}
\begin{CornellChecklist}
\item \textbf{Scoping}
Record real, effective, saved, and supplementary identities at each sensitive operation.
\item \textbf{Attack-surface identification}
Minimize privileged code and inventory every input accepted before privilege is reduced.
\item \textbf{Input and data-flow analysis}
Check every identity-changing return value and test attempted regain.
\item \textbf{Trust-boundary review}
Bracket the minimal operation and verify state on normal and exceptional exits.
\item \textbf{Candidate-point inspection}
Audit permissions and ownership on the full path, not only the final object.
\item \textbf{Control-flow review}
Use atomic creation with restrictive initial permissions and verify the opened object.
\item \textbf{Resource and lifetime review}
Avoid privileged resolution through untrusted directories and use trusted directory handles where possible.
\item \textbf{Error-path analysis}
Bind policy to a resolved object or trusted directory-relative operation rather than a string alone.
\item \textbf{Mitigation validation}
Verify type, owner, permissions, and flags on the opened descriptor before use.
\item \textbf{Evidence and reporting}
Use no-follow and directory-relative primitives where supported, then verify the opened object.
\item \textbf{Scoping}
Review link-creation permissions and avoid treating pathname uniqueness as object identity.
\item \textbf{Attack-surface identification}
Replace check-then-use sequences with one atomic operation or descriptor-based validation.
\end{CornellChecklist}
\section{Chapter Synthesis}
UNIX file security is fundamentally about identity, namespace control, and atomicity. A reviewer must follow effective authority, inspect every directory in a path, prefer stable handles over repeated names, and eliminate check-use gaps. Prioritize privileged operations in writable directories, unsafe temporary-file creation, incomplete privilege drops, and descriptor or stream errors that alter protected data.
\section{Self-Test Questions}
\begin{enumerate}
\item Which identity does the kernel use for an access decision?
\item Why are set-user-ID and similar programs high risk?
\item What proves that privilege cannot be regained?
\item What additional state must be controlled?
\item What authority follows from directory permissions?
\item How should a new file acquire secure permissions?
\item \textbf{Scenario:} A privileged program checks that a temporary pathname is unused and then opens it for writing. Why is this unsafe?
\item \textbf{Scenario:} A process changes only its effective user ID before parsing untrusted input and retains a saved privileged ID. Has it permanently dropped privilege?
\end{enumerate}
\clearpage
\subsection*{Answer Key}
\begin{enumerate}
\item Processes carry multiple user and group identities, and files carry owners and mode bits. The relevant identity can differ across system calls and privilege transitions.
\item They process less-trusted input while holding another principal's authority. Environment, files, descriptors, signals, and invoked programs become part of the attack surface.
\item A permanent drop must update all relevant identities and groups, check errors, and verify the resulting state. Partial changes may leave a saved identity or capability that restores authority.
\item Temporarily changing effective identity can support a narrow task, but restoration, signals, threads, descriptors, and concurrent work complicate the model. Untrusted work must not run while elevated.
\item Directory search, create, remove, and rename rights differ from permissions on the target file. A protected file can still be replaced or redirected when an attacker controls a containing directory.
\item Creation mode, process mask, inherited context, flags, and preexisting names determine the resulting object. Opening and then fixing permissions leaves a window and may operate on an attacker-chosen file.
\item An attacker can create or redirect the name between the check and open. Use an atomic exclusive-creation primitive in a trustworthy directory and operate through the returned descriptor.
\item No. The saved identity may permit regain, and parsing occurs within a complex reversible state. A permanent drop must update all relevant credentials, supplementary groups, and verify that regain fails.
\end{enumerate}
\section{Key-Term Recap}
\begin{longtable}{>{\raggedright\arraybackslash}p{0.27\textwidth}p{0.65\textwidth}}
\toprule
\textbf{Term} & \textbf{Working meaning for review} \\
\midrule
\endfirsthead
\toprule
\textbf{Term} & \textbf{Working meaning for review} \\
\midrule
\endhead
Users, groups, and processes & Processes carry multiple user and group identities, and files carry owners and mode bits. \\
Privileged programs & They process less-trusted input while holding another principal's authority. \\
Permanent privilege drop & A permanent drop must update all relevant identities and groups, check errors, and verify the resulting state. \\
Temporary privilege drop & Temporarily changing effective identity can support a narrow task, but restoration, signals, threads, descriptors, and concurrent work complicate the model. \\
File and directory permissions & Directory search, create, remove, and rename rights differ from permissions on the target file. \\
File creation & Creation mode, process mask, inherited context, flags, and preexisting names determine the resulting object. \\
Directory safety and dangerous places & World-writable or attacker-controlled directories permit name replacement, link insertion, and cleanup races. \\
Filenames and paths & Repeated separators, parent components, alternate roots, mount changes, and links can make textual form diverge from the resolved object. \\
Descriptors, inodes, and directories & A descriptor refers to an opened object independently of later renaming, while an inode identifies filesystem metadata within its context. \\
Symbolic links & A symbolic link stores another path and can redirect lookup across directories or filesystems. \\
\bottomrule
\end{longtable}
\section{One-Sentence Takeaway}
\begin{CornellExamBox}{One-Sentence Takeaway}
\textbf{Secure UNIX resource handling binds least privilege to atomic, descriptor-based operations on verified objects rather than mutable pathnames.}
\end{CornellExamBox}
\end{document}
